\documentclass[10pt,a4paper]{article}
\usepackage[margin=23mm]{geometry}
\usepackage{amsmath,amssymb,booktabs,array,longtable}
\usepackage[hidelinks]{hyperref}
\setlength{\parindent}{0pt}
\setlength{\parskip}{4pt}
\setlength{\emergencystretch}{2em}
\newcommand{\dd}{\mathrm d}
\newcommand{\pos}[1]{[#1]_+}
\title{Urban service allocation:\\traffic and sensitivity specification}
\author{Urban traffic methods}
\date{28 October 2026}
\usepackage{xcolor}
\definecolor{revisionred}{HTML}{D13F43}
\newcommand{\rev}[1]{{\color{revisionred}#1}}
\begin{document}
\maketitle

\textbf{Purpose and scope.}
This document specifies a discrete continuous-mass model for an urban
signal-control task. It defines lane admission, order-dependent access,
service and complete population-time sensitivity. The representative
approach has two lanes upstream of a finite left-turn bay and three lanes
at the stop line. Numerical calibration, baseline implementation checks
and empirical validation are separate from this mathematical definition.

\section{Traffic world, action and physical ownership}

Retain one isolated four-leg junction, fixed routes and four protected
effective-service groups: NS through/right, NS left, EW through/right,
EW left. With physical capacity $C_\ell$ and saturation $s_{\ell m}$,
\rev{the service law below adapts Han et al.'s average-signal interpretation to the task's four groups and capacity bound.}
Its connection to the source is explained in Section 5.
\begin{equation}
 g_p\geq0,\quad \sum_{p=1}^4g_p=1,\qquad
 \mu_{\ell m}(g)=\min\{C_\ell,s_{\ell m}g_{p(m)}\}.
 \label{eq:action}
\end{equation}
Mass is continuous vehicle-equivalent quantity; $C,s,\mu$ are veh-eq/s.
A feasible direction preserves the simplex. Route labels, permissions,
topology and source-order skeleton are structural. No phase waveform,
integerization, microscopic lane changing or endogenous route choice is
introduced. Control uses constant $g$ during the declared control window,
then a common candidate-independent clearance allocation.

The representative approach has two parallel feeder lanes and three
stop-line lanes.
\rev{The following graph defines the representative task geometry:}
\begin{equation}
 U\longrightarrow\{F_1,F_2\},\qquad
 F_1\longrightarrow\{B,T_1\},\qquad F_2\longrightarrow T_2 .
 \label{eq:geometry}
\end{equation}
$U$ is the aggregate upstream road. Within the local zone, $F_1$ carries
left/through traffic and $F_2$ through/right traffic. The internal bay
opening divides $F_1$ into left-only bay $B$ and through-only continuation
$T_1$; the same cross-section divides $F_2$ from shared through/right
continuation $T_2$. In the symmetric core
$L_{F_1}=L_{F_2}$ and $L_B=L_{T_1}=L_{T_2}$, giving equal total approach
path length. All sections own disjoint physical road space and storage.

Left traffic may enter only $F_1$, then $B$; right traffic only $F_2$,
then $T_2$. Through traffic may choose $F_1$ or $F_2$ at actual admission
from $U$, but follows its chosen lane thereafter: $F_1$ to $T_1$ or
$F_2$ to $T_2$. There are no cross-feeder lane-change edges and no direct
$U$ to $B,T_1,T_2$ edges. The queue estimate at the choice boundary uses
current receiving-feeder mass and earlier actual admissions, as in the
allocation rule below; no downstream look-ahead or new driver cost is added.
The two physical lanes can therefore differ in accessibility without
allowing traffic to escape a lane it already occupies.

V0 uses movement-exclusive full-length paths. V1 uses a shared lane.
V2/V3 share the representative bay geometry and permissions: V2 suppresses
order-sensitive obstruction by construction, whereas V3 tests it.

V4 reserves demand/turn-share, bay-length/storage, arrival-order and
downstream-receiving combinations, including asymmetric approaches,
from calibration and method tuning. Those scenario variations test
robustness after the isolated and combined mechanism ablations.
\par The smaller single-feeder
case remains a local mechanism diagnostic. Permissions may differ
between scenario families, but not between compared models or paired
order cases of the same scenario.

\textbf{Processing order and physical lane order.}
An unassigned link-end or external-source pool has an ordered processing
record. A realized local lane has an actual first-in, first-out sequence.
The input record governs allocation inspections; the occupied lane also
imposes front blocking. A waiting head in $F_1$ cannot be passed by a
later resident, while an independent $F_2$ has its own available path.
A record cannot be relabeled as unassigned to permit physical overtaking.

\section{Propagation and the step clock}

Each physical section retains cumulative counts $N^+_{e,k},N^-_{e,k}$
and occupancy $n_{e,k}=N^+_{e,k}-N^-_{e,k}$ at $t_k=k\Delta t$.
Use triangular-FD capacity $C_e$, free-flow delay $\tau^f_e$, backward
delay $\tau^b_e$ and storage $K_e=C_e(\tau^f_e+\tau^b_e)$, with
$\Delta t\leq\min_e(\tau^f_e,\tau^b_e)$.
\rev{The storage identity and timestep domain are inherited from the project's triangular-FD LTM transcription.}
Road capacities, delays and legal saturation rates are finite and positive.
An incompatible short section requires a compatible discretization;
no fictitious delay, clipping or smoothing is inserted.

\rev{Reuse that transcription's affine boundary-history reconstruction.}
With empty prehistory and $\theta=(t-t_h)/\Delta t$,
\begin{equation}
 {\cal H}^{\pm}_e(t)
 =(1-\theta)N^\pm_{e,h}+\theta N^\pm_{e,h+1},
 \qquad t_h\leq t\leq t_{h+1}.
 \label{eq:history}
\end{equation}

The propagation basis is Yperman's LTM, using the repository's verified
\href{../../../../docs/milestones/milestone_0_ltm_baseline.md}{discrete transcription}.
The \href{https://www.researchgate.net/publication/28360292_The_Link_Transmission_Model_for_dynamic_network_loading}{author-uploaded thesis}
is the literature source; the internal feeder graph is the present
physical composition.
\par

The sending count subtracts prior departures from entries early enough
to reach the exit. The receiving count starts with storage, adds departures
whose space-release wave has reached the entrance, and subtracts all
previous entries. Both counts are bounded by capacity times timestep.
\rev{The two limits below reuse that discrete LTM transcription on the task's internal section graph:}
\begin{align}
 A_{e,k}&=\pos{{\cal H}^+_e(t_{k+1}-\tau^f_e)-N^-_{e,k}},
 &S_{e,k}&=\min(C_e\Delta t,A_{e,k}),\label{eq:sending}\\
 R_{e,k}&=\min\{C_e\Delta t,\,
 \pos{K_e+{\cal H}^-_e(t_{k+1}-\tau^b_e)-N^+_{e,k}}\}.
 \label{eq:receiving}
\end{align}
The positive part sets negative availability to zero. Free-flow delay
controls downstream travel; backward delay controls how soon new room
is available upstream. $A$ is ready resident mass for this update;
$S,R$ are mass allowances,
not instantaneous rates. Queries use already stored boundaries.
New inflow cannot traverse its receiving section in the same update.

This is a timestep model. The internal head-service calculation produces
endpoint totals; it does not add propagation-history knots.
Readiness uses the inherited discrete sending convention, rather than
claiming every selected resident was ready at every instant in the step.
No event-history refinement is adopted.

These equations and the left-interval objective are inherited from the
approved repository M0 transcription, checked against the local LTM and
objective code. Their application to the internal section graph is the
declared urban composition, not existing implemented M2 behavior.

\section{Ledger, eligibility and lane membership}


Link-wise route tracking on a macroscopic backbone is a literature idea;
\href{https://doi.org/10.1016/j.simpat.2025.103088}{de Souza et al. (2025)}
achieve individual tracking using integer flows.
\rev{That reference motivates identity tracking; no integer-flow formula is imported here.}
The representation here
instead retains continuous packet quantities with fixed route/order
metadata. A long link has a route-mass ledger; each realized local lane
has its own ordered ledger. The continuously changing mass and timing
fields, rather than route IDs or permutation choices, carry derivatives.
The exact split/merge and timing construction follows the repository's
\href{../../../../docs/milestones/milestone_1_ordered_route_mass_ledger.md}{continuous route-mass specification}.
\par

Store positive-mass cohorts, removing exact zero entries without a cutoff.
A cohort stores route-table reference, route position, continuous mass,
current owner and affine eligibility provenance. Current movement is
projected from its route. Internal $U/F_1/F_2/B/T_1/T_2$ crossings preserve route
position; the actual junction crossing advances it. Future movement
sequences are not duplicated in grid entries.

Only actual admission establishes receiving-lane membership. Unadmitted
mass stays with its old owner without a binding future-lane vector.
Traffic already in $F_1$ or $F_2$ retains that realized lane; its movement
and the declared connectivity determine its continuation. Resident traffic
cannot be reassigned across feeders or terminal lanes to improve balance.

Retain the exact source order of the unserved record. Before evaluating
an operator, coalesce exactly compatible adjacent representation fragments:
same route/position, owner, permissions and contiguous affine eligibility
with the M1 exact line condition. No tolerance-based mixing or route
merging is used. Serialization changes must preserve predictions; they
do not define a new physical traffic unit.


The near-intersection movement grid and the lane ledger represent the
same mass in complementary forms. The grid supplies per-lane continuous
movement amounts and loads; the ledger resolves the exact order hidden
by a mixed bin. Classical
\href{https://doi.org/10.1016/0191-2615(94)90002-7}{CTM (Daganzo, 1994)}
uses physical-space cells. Here fixed cells are cumulative-mass/rank
intervals: a mass-coordinate adaptation, not an additional spatial CTM
propagation law. Physical capacities, travel and storage remain governed
by the section LTM histories and geometry.

Use a fixed equal-mass lattice with geometry-derived capacity masks,
including a possible partial final bin. For packet $r$ on lane $\ell$,
let $I_{\ell r}=[z^-_{\ell r},z^+_{\ell r}]$, where
$z^-_{\ell r}=\sum_{v<r}q_{\ell v}$ and
$z^+_{\ell r}=z^-_{\ell r}+q_{\ell r}$.
Bin $c$ has fixed boundaries $\zeta_{\ell c},\zeta_{\ell,c+1}$.
\rev{The overlap formula is derived from interval intersections in the project's fixed cumulative-mass grid; CTM supplies cell-based inspiration, not this flux-free representation:}
\begin{equation}
 \begin{aligned}
 \omega_{\ell cr}&=
 [\min(z^+_{\ell r},\zeta_{\ell,c+1})
           -\max(z^-_{\ell r},\zeta_{\ell c})]_+,\\
 M_{\ell cm}&=\sum_{r:m_r=m}\omega_{\ell cr},\qquad
 n_\ell=\sum_{c,m}M_{\ell cm}=\sum_rq_{\ell r}.
 \end{aligned}
 \label{eq:grid}
\end{equation}
Here $\omega_{\ell cr}$ is packet $r$'s overlap mass in bin $c$,
$M_{\ell cm}$ is movement $m$'s mass in that lane/bin, and $n_\ell$
is total resident lane mass, including moving and waiting traffic.
Their units are vehicle-equivalents, not metres.
This formula reconstructs the grid from the ledger; it is not Daganzo's
cell flux equation. Rebuild it atomically after conservative service and
admission, keeping zero cells in fixed padded shapes. Changing the mass
resolution changes neither storage nor geometry. The grid is a derived
continuous state, not another stock to add to the mass account. A mixed
bin cannot reconstruct the head sequence, so exact lane access/service
continues to use the ledger.
\par

\textbf{Sending limits transfers, not candidate inspections.}
Form candidates from ready mass $A$, cutting an eligibility segment
exactly if needed. For an unassigned pool, scan these candidates and
spend at most $S$ in total. The candidate list is not restricted to the first $S$ units before access
is checked: an earlier blocked movement must not hide later ready traffic
with an independently accessible path. For a physical lane, the access rule below forces
actual served mass to be a prefix.

This is an explicit adaptation of M1's prefix extraction. After a
selection that skips an independently blocked movement, unserved mass
keeps its original entry/travel history and relative residual order.
Waiting for admission does not restart its travel delay. Every past departure was eligible when served, and eligibility
is monotone. Total eligible residual mass therefore remains
${\cal H}^+(t_{k+1}-\tau^f)-N^-_k$ even if eligible departures were not
a global processing-order prefix. The adaptation is not claimed to exist
in current M1 code, nor to establish passing on a represented single lane.

\section{Allocation, unassigned residuals and physical access}


The equilibrium allocation idea is motivated by
\href{https://doi.org/10.1016/j.trb.2025.103208}{Gong, Bliemer and Raadsen (2025)}.
The raw queue-mass cost and capped per-packet entry rule below are the
present adaptation. They are derived from a strictly convex load-balancing
problem, not claimed as that paper's original exit-resistance equation
or general algorithm.
\par

Let $a$ be the sending owner, $r$ a ready cohort of mass $q_r$ and
movement $m_r$, and $\ell$ a receiving lane/section.
Permission $P_{m\ell}\in\{0,1\}$ and reachability
$E_{r\ell}\in\{0,1\}$ encode actual paths/access, not a future green plan.

Initialize residual sending $s=S_{a,k}$, residual receiving
$R^{\rm rem}_\ell=R_{\ell,k}$ and perceived loads $h_\ell=n_{\ell,k}$.
Process cohorts in recorded order. Let ${\cal B}$ contain movements
already left partly unserved in this scan. Each sending owner has its own
scan and blocked-movement set; one inbound leg does not block another.
If $m_r\in{\cal B}$ set $V_r=x_{r\ell}=0$; otherwise admit the largest
feasible total $V_r$, limited by packet mass, remaining sending and
legal reachable receiving. Distribute that total to reduce differences
in current lane loads, using a common filling level $\lambda_r$.
\par\begin{samepage}
\rev{The quadratic objective and its capped solution below are defined and derived for this entry problem; Gong et al. supply the equilibrium motivation.}

\begin{align}
 u_{r\ell}&=P_{m_r\ell}E_{r\ell}R^{\rm rem}_\ell,\qquad
 V_r=\min\{q_r,s,\sum_\ell u_{r\ell}\},\nonumber\\
 x_r&=\underset{0\leq x_{r\ell}\leq u_{r\ell},\
                    \sum_\ell x_{r\ell}=V_r}{\arg\min}
          \frac12\sum_{\ell:P_{m_r\ell}=1}(h_\ell+x_{r\ell})^2,
 \nonumber\\
 x_{r\ell}&=\min\{u_{r\ell},\pos{\lambda_r-h_\ell}\},\qquad
                    \sum_\ell x_{r\ell}=V_r .
 \label{eq:allocation}
\end{align}
\par\end{samepage}

The solver chooses $\lambda_r$ so the lane additions sum to $V_r$.
Its units are vehicle-equivalent mass: $\lambda_r-h_\ell$ is the mass
needed to bring a lane to that load. The positive part excludes lanes
already above the level and $u_{r\ell}$ caps each feasible admission.
This is a mass-balancing rule, not a stop-line service cap. Illegal lanes
have zero addition. Strict convexity selects unique additions; the filling
level itself need not be unique when all usable lanes are capped.
For $V=0$, use $x=0$. Commit only actual entry.
\rev{The following resource updates are conservation identities derived from the admitted amounts in \eqref{eq:allocation}:}
\begin{equation}
 s\leftarrow s-V_r,\quad
 R^{\rm rem}_\ell\leftarrow R^{\rm rem}_\ell-x_{r\ell},\quad
 h_\ell\leftarrow h_\ell+x_{r\ell},\quad
 q_r^{\rm rem}=q_r-V_r .
 \label{eq:allocation-update}
\end{equation}
If $q_r^{\rm rem}>0$, retain it at its original position and add $m_r$
to ${\cal B}$ for this update. Later cohorts of that movement wait;
other movements are considered. Reset ${\cal B}$ next update.
The retained record then puts the older residual first within its
movement. There is no waiting-target allocation or old-target rationing law.

$h$ is the selected queue-perception estimate: pre-state local lane mass
plus earlier actual admissions. It contains no invented waiting targets
and does not subtract forecast signal discharge. Actual next occupancy
separately includes that step's exits. Raw queue mass is the first
perception model; no service-aware cost or unequal preference is added.

\subsection*{Access in the admitted layouts}

For an unassigned pool with independently accessible receiving lanes,
$E_{r\ell}=1$ on those paths. Movement labels alone do not create access.
This applies at $U\to\{F_1,F_2\}$ on a packet's legal feeder choices,
and to a source feeding independent V0 paths. Once mass occupies $F_1$
or $F_2$, it cannot use the other feeder's capacity.

For an occupied single lane, including $F_1$ at its fork,
\rev{the access predicate is defined from realized-lane FIFO and the chosen connectivity:}
\begin{equation}
 E_{r\ell}
 =\mathbf1\{\text{every earlier resident cohort has fully crossed}\}
 \quad\text{on a legal outgoing edge}.
 \label{eq:physical-access}
\end{equation}
An older physical residual prevents later physical entry even while the
allocator inspects later records. Several actual lanes have separate
ledgers/access conditions, not one global FIFO.

Movement-specific blocking governs the unassigned scan; actual lane
occupation adds the physical access constraint. At $U$, a full $F_1$
prevents left entry while through traffic may still use reachable $F_2$.
At the fork, a through head in $F_1$ can obstruct a later turner despite
an empty bay. A full bay can similarly leave a left head blocking later
through traffic in $F_1$. Independently resident $F_2$ traffic may continue
toward $T_2$ if its own sending, receiving and order permit it.

A zero left-turn stop-line rate is not automatically zero lane-entry room.
A left turner may join a reachable feeder or bay with room and subsequently
contribute to spillback. In this geometry it cannot turn left from $T_1$
or use $F_2,T_2$; a full bay has no alternative left-turn stop-line lane.
Each constraint refers to its actual physical boundary.

\subsection*{Direct definition checks}

The first case below is a local operator diagnostic, not a new primary
scenario or a claim that every aggregate road offers independent access.
For an independent-access pool with order $[L(0.4),T(0.6)]$, $A=1$,
$S=0.6$ and exclusive-destination receiving $(0,0.5)$, left entry is
zero and through entry $0.5$. Residual order is $[L(0.4),T(0.1)]$.
Truncating candidates to the first $S$ mass would expose only $0.2$
through mass and give the wrong result for this pool rule.

For actual $F_1$ take $[T(0.4),L(0.6)]$, $S_{F_1}=1$, $R_{T_1}=0.2$,
$R_B=1$. Through entry is $0.2$, left entry zero.
Reverse the order with the same totals and branch loads
$(n_B,n_{T_1})=(0,1.8)$. Left uses $B$ first, then through enters $0.2$:
total $0.8$ rather than $0.2$.
These branch bounds are compatible with
$\Delta t=\tau^f=\tau^b=1$, $C_B=C_{T_1}=1$, $K_B=K_{T_1}=2$.
If an independent ready $0.3$ in $F_2$ has enough receiving in $T_2$,
it enters in either case. Cross-section totals are then $0.5$ and $1.1$.
These are prepared hand states, not a direct initializer or rollout.

At the $U$ choice boundary, a compatible balancing case has feeder loads
$(1.2,1.0)$, receiving $(0.8,1.0)$, ready through mass $1$ and sending $0.6$.
Additions $(0.2,0.4)$ leave perceived loads $(1.4,1.4)$;
the remaining upstream $0.4$ is unassigned.

\section{Elapsed service and selected receiving rates}


The use of green fractions as average service follows the continuum
signal interpretation of
\href{https://arxiv.org/abs/1309.3461}{Han et al. (2014)}.
The conservation/supply framework follows first-order node modeling;
\href{https://doi.org/10.1016/j.trb.2010.06.004}{Tampere et al. (2011)}
is a basis for those requirements. The particular green-weighted leg
sharing and equal-within-leg capped reuse below are model-specific
priorities, rather than a copy of the reference's oriented-capacity rule.
\par

A stop-line lane exposes its current ready resident head, with sending
quota $S_\ell$ and cap $c_\ell=\mu_{\ell m}(g)$.
\rev{The constant-rate conversion below is the project's chosen average-receiving convention, applied to the LTM allowance:}
\begin{equation}
 \rho_j=R_{j,k}/\Delta t .
 \label{eq:receiving-rate}
\end{equation}
This is a timestep computational average, not an inferred within-step
physical supply curve.

For each output $j$, collect current heads by inbound leg $i$.
A leg-to-output movement here has one service group.
\rev{The next two equations define this model's green-weighted leg and equal-lane priority rules within the cited feasibility framework:}
\begin{align}
 C_{ij}&=\sum_{\ell\in{\cal L}_{ij}}c_\ell,\quad
 w_{ij}=g_{p(i,j)},\quad W_j=\min(\rho_j,\sum_iC_{ij}),\nonumber\\
 r_{ij}&=\min(C_{ij},\theta_jw_{ij}),\quad
                    \sum_i r_{ij}=W_j,\label{eq:leg-sharing}\\
 r_\ell&=\min(c_\ell,\eta_{ij}),\quad
                    \sum_{\ell\in{\cal L}_{ij}}r_\ell=r_{ij}.
 \label{eq:lane-sharing}
\end{align}
Remove zero-cap competitors; zero $W$ gives zero rates.

Here $\theta_j$ is chosen so green-weighted leg rates total $W_j$;
$\eta_{ij}$ is the common uncapped lane rate chosen so the lane rates
total $r_{ij}$. Its units are vehicle-equivalents/s, unlike the mass-level
$\lambda_r$ used for entry. Each lane receives its cap or that common
rate, whichever is smaller. Thus a grant $0.6$ splits $0.3/0.3$ between
two unrestricted lanes; with one capped at $0.1$, the other may use $0.5$.
The leg weight is counted once, independent of lane count.
\par


Rates describe the current heads; their one-step amounts require the
time needed to clear each head. For a one-second shared-lane example,
a ready front block of mass $0.2$ at rate $0.4$ uses half a second,
leaving half a second for an already-ready following movement at rate
$0.6$, which then serves $0.3$. With several lanes, a changed head can
alter receiving competition. Let $\sigma$ be time already accounted for,
starting at zero, and $q_\ell$ the residual head mass. Current rates apply
until a head clears, remaining sending is spent or the step ends:
\par
\rev{Combining the current-head, sending and timestep stopping constraints derives the duration rule:}
\begin{equation}
 \delta=\min\left\{\Delta t-\sigma,\
 \min_{\ell:r_\ell>0}
       \frac{\min(q_\ell,S_\ell^{\rm rem})}{r_\ell}\right\}.
 \label{eq:head-time}
\end{equation}
Transfer $r_\ell\delta$, process tied events together, expose new heads
and recompute sharing. Stop when rates vanish or time ends.
For a chosen lane/movement stream, $r_a$ is its rate in interval $a$,
zero when it is not the current head.
\rev{Summing the transferred mass rate times duration derives its recorded timestep amount:}
\begin{equation}
 y=\sum_a r_a\delta_a.
 \label{eq:service-amount}
\end{equation}
Include remaining idle time with zero rate when describing the full
timestep. This accounting does not add propagation-history knots.
A zero-rate head blocks its lane. Newly admitted mass is outside this
resident service tranche.

A later ready head joins current competition; earlier discharge is not
compensated to equalize final totals. Constant receiving removes a
separate consumable-pool exhaustion rule, but head clearance and changing
competitors remain sequential. Receiving reuse does not transfer one
signal group's service to another.

For V1's separate shared left/through service diagnostic, ready left
mass $q_L$ precedes sufficient through mass. With nonbinding receiving
and physical caps, its clearance time determines the remaining through
service time.
\rev{The example below is derived by applying \eqref{eq:service-amount} to a head-clearing left block and its through follower:}
\begin{equation}
 y_L=q_L,\qquad
 y_T=s_Tg_T\left(\Delta t-\frac{q_L}{s_Lg_L}\right),
 \quad 0<\frac{q_L}{s_Lg_L}<\Delta t.
 \label{eq:head-example}
\end{equation}
With unit timestep/saturations, $q_L=0.2$ and $(g_L,g_T)=(0.4,0.6)$,
through discharge is $0.3$. This is the optional algebraic check of the
rate-versus-time mechanism explained in the main report. For
$g_L(\alpha)=g_L+\alpha$ and $g_T(\alpha)=g_T-\alpha$, differentiating
\eqref{eq:head-example} \rev{by the product and quotient rules gives}
\begin{equation}
 \frac{\mathrm dy_T}{\mathrm d\alpha}
 =-s_T\left(\Delta t-\frac{q_L}{s_Lg_L}\right)
    +\frac{s_Tg_Tq_L}{s_Lg_L^2}.
 \label{eq:head-example-derivative}
\end{equation}
At the stated values the rate-loss term is $-0.5$ and the earlier-clearance
term is $0.75$, giving $0.25$. These are the two channels of
\eqref{eq:service-derivative}; \eqref{eq:duration-derivative} supplies
the head-time response. A fully cleared fixed head contributes no extra
mass itself; its followers can use the time saved. This is not a
complete-TST benefit or an aggregate-rival failure claim.

\section{Paired transfers, histories and feasibility}

A one-in/one-out internal edge transfers its ready prefix up to
$\min(S_{\rm up},R_{\rm down})$. Forks use allocation/access; the central
junction uses elapsed service. An external source adds $D_{o,k}$ to its
waiting record before admission, with budget $Q_{o,k}+D_{o,k}$.
A sink removes at most sending and its declared allowance.

Each internal entry interface has one sending owner; $U$ has two
receiving feeder choices sharing its total sending allowance. Shared receiving at the
central signal node uses the simultaneous current-head sharing law, rather
than an arbitrary serial order of inbound legs.
All nodes read pre-state histories/ledgers. Each section's sending and
receiving is spent once at its boundary. New admissions are appended
after selecting resident service; today's exits create no immediate
receiving credit.
\rev{The cumulative-count update below reuses the project's conservative LTM/ledger bookkeeping:}
\begin{equation}
 N^+_{e,k+1}=N^+_{e,k}+b_{e,k},\quad
 N^-_{e,k+1}=N^-_{e,k}+y_{e,k},\quad
 n_{e,k+1}=n_{e,k}+b_{e,k}-y_{e,k}.
 \label{eq:state-update}
\end{equation}
Internal flow leaves one owner and enters the next with identical mass
and route identity. Source admission removes source waiting; sink exit
removes road mass. No residual is duplicated as lane stock.


The next update needs to distinguish traffic that has reached service
from mass still travelling. Each admitted packet therefore carries an
entry interval consistent with the uniformly reconstructed cumulative
history. In a receiving section with total entrance $b>0$, a child of
mass $q^{\rm in}$ preceded by admitted mass $z$ receives the interval below.
\rev{Its endpoints are derived from the uniform boundary reconstruction inherited from the continuous ledger specification:}
\begin{equation}
 [t^{\rm in}_{\rm front},t^{\rm in}_{\rm tail}]
 =\left[t_k+\Delta t\,\frac z b,\
       t_k+\Delta t\,\frac{z+q^{\rm in}}b\right].
 \label{eq:entry-time}
\end{equation}
Both endpoints are entry times: the first and last parts of that child.
For $b=1$, $z=0.2$, $q^{\rm in}=0.3$ and unit timestep, the child enters
from $0.2$ to $0.5$ seconds after the update starts. Add section free-flow
travel time $\tau^f_e$ to both endpoints to obtain its earliest exit
readiness interval. Retained portions keep that timing; waiting does not
restart travel. A later timestep can select only the ready fraction.
Zero inflow has no admitted child and therefore no entry interval.

Children are appended in actual admission order. For cross-input
outgoing bookkeeping, order by service-interval start, then fixed input
leg, lane and source-cohort indices. That serialization does not change
shared-rate grants. The outgoing roads lead to sinks, so it cannot create
a later signal-order advantage in this task.
\par

The timestep domain gives $S_{e,k}\leq n_{e,k}$ and
$R_{e,k}\leq K_e-n_{e,k}$.
\rev{Combining these inherited bounds with \eqref{eq:state-update} derives the storage implication below:}
\begin{equation}
 0\leq b_{e,k}\leq R_{e,k},\quad0\leq y_{e,k}\leq S_{e,k}
 \ \Longrightarrow\
 0\leq n_{e,k}+b_{e,k}-y_{e,k}\leq K_e .
 \label{eq:storage-bound}
\end{equation}
\rev{Summing the paired conservation updates from empty loading, including source waiting and sink exits, derives}
\begin{equation}
 \sum_oQ_{o,k}+\sum_en_{e,k}+Z_k=D_{<k},
 \label{eq:mass-identity}
\end{equation}
also route by route. $Z_k$ is cumulative sink exit; $D_{<k}$ all earlier
exogenous release.

A finite state has a defined unique result: finite source scan, unique
convex allocation vectors, unique capped-sharing vectors and a finite
head-event calculation. Each positive service interval exhausts a
head/quota or reaches step end. Exact ties are handled together; no
mass cutoff or software priority manufactures uniqueness. These are
definition-level arguments, not implementation or traffic-fidelity evidence.


\section{Sensitivities through the traffic modules}
Let $g(\alpha)=g+\alpha d$, with $\sum_p d_p=0$, define a feasible
service transfer. All derivatives below are with respect to that same
control amount $\alpha$. Fix routes, permissions and the current active
head/constraint pattern. A structural switch requires its own directional
branch.

The path is service fractions to head rates, then actual transfers, then
continuous ledger masses/timing, grid overlaps and road histories, then
later sending/receiving and admission, and finally sink exits and complete
time. The grid is reconstructed inside the step map; labels and order are
structural while masses and timing are differentiated.

\par\begin{samepage}
Differentiate the admission rule and its mass balance in
\eqref{eq:allocation}. For the current packet, let ${\cal A}$ contain
positive uncapped lanes ($0<x_\ell<u_\ell$) and ${\cal U}$ lanes on
binding upper bounds. Recall $V$ is total admitted mass, $h_\ell$ the
load before this allocation, $u_\ell$ its entry bound and $\lambda$
the common filling level.
\rev{The balance and cases below are derivatives of \eqref{eq:allocation} on fixed active categories:}
\begin{equation}
 \begin{aligned}
 \frac{\mathrm dV}{\mathrm d\alpha}
 &=\sum_\ell\frac{\mathrm dx_\ell}{\mathrm d\alpha},\\
 \frac{\mathrm dx_\ell}{\mathrm d\alpha}
 &=\begin{cases}
 \displaystyle\frac{\mathrm d\lambda}{\mathrm d\alpha}
             -\frac{\mathrm dh_\ell}{\mathrm d\alpha},&\ell\in{\cal A},\\
 \displaystyle\frac{\mathrm du_\ell}{\mathrm d\alpha},&\ell\in{\cal U},\\
 0,&\text{no addition on the unchanged lower branch}.
 \end{cases}
 \end{aligned}
 \label{eq:allocation-derivative}
\end{equation}
\par\end{samepage}
Substituting these lane responses into the total and solving for the
common filling-level sensitivity \rev{in \eqref{eq:allocation-derivative} derives the algebraic solution}
\begin{equation}
 \frac{\mathrm d\lambda}{\mathrm d\alpha}
 =\frac1{|{\cal A}|}\left(
 \frac{\mathrm dV}{\mathrm d\alpha}
 -\sum_{\ell\in{\cal U}}\frac{\mathrm du_\ell}{\mathrm d\alpha}
 +\sum_{\ell\in{\cal A}}\frac{\mathrm dh_\ell}{\mathrm d\alpha}\right).
 \label{eq:allocation-level-derivative}
\end{equation}
The capped-bound response consumes part of the admitted-mass response;
starting-load increases require a higher level to preserve free additions.
The denominator counts positive uncapped receivers. Evaluate this level
response first, then the individual additions. This carries prior traffic
changes into actual later lane entry, not a new direct dependence on green.

Stable fully capped branches follow their caps; stable zero admission
has zero admission derivative. Differentiate $V$ through its limiting
argument and $u$ through the current receiving/access branch. Sequential
updates differentiate the resource updates in \eqref{eq:allocation-update}
and preserve the response of remaining mass and resources.
\rev{The following is direct differentiation of those resource identities, not a new allocation rule:}
\begin{equation}
 \begin{aligned}
 \frac{\mathrm dh_\ell}{\mathrm d\alpha}&\leftarrow
  \frac{\mathrm dh_\ell}{\mathrm d\alpha}
                  +\frac{\mathrm dx_\ell}{\mathrm d\alpha},\\
 \frac{\mathrm dR_\ell^{\rm rem}}{\mathrm d\alpha}&\leftarrow
  \frac{\mathrm dR_\ell^{\rm rem}}{\mathrm d\alpha}
                  -\frac{\mathrm dx_\ell}{\mathrm d\alpha},\\
 \frac{\mathrm ds}{\mathrm d\alpha}&\leftarrow
  \frac{\mathrm ds}{\mathrm d\alpha}-\frac{\mathrm dV}{\mathrm d\alpha},
 \qquad
 \frac{\mathrm dq^{\rm rem}}{\mathrm d\alpha}
  =\frac{\mathrm dq}{\mathrm d\alpha}-\frac{\mathrm dV}{\mathrm d\alpha}.
 \end{aligned}
 \label{eq:scan-derivative}
\end{equation}
Green directly changes the signal-limited rate in \eqref{eq:action},
and changes the weights in \eqref{eq:leg-sharing}; receiving changes
through \eqref{eq:receiving-rate}. Recall $r_{ij}$ is the leg grant,
$C_{ij}$ its cap, $w_{ij}$ its weight, $W_j$ the usable total and
$\theta_j$ the common multiplier. For uncapped legs ${\cal F}$ and
capped legs ${\cal C}$,
\rev{differentiating the model-defined \eqref{eq:leg-sharing} on fixed cap categories derives}
\begin{equation}
 \begin{aligned}
 \frac{\mathrm dW_j}{\mathrm d\alpha}
 &=\sum_i\frac{\mathrm dr_{ij}}{\mathrm d\alpha},\\
 \frac{\mathrm dr_{ij}}{\mathrm d\alpha}
 &=\begin{cases}
 \displaystyle\frac{\mathrm dC_{ij}}{\mathrm d\alpha},&i\in{\cal C},\\
 \displaystyle w_{ij}\frac{\mathrm d\theta_j}{\mathrm d\alpha}
       +\theta_j\frac{\mathrm dw_{ij}}{\mathrm d\alpha},&i\in{\cal F}.
 \end{cases}
 \end{aligned}
 \label{eq:sharing-derivative}
\end{equation}
Substitute these responses into the total-rate balance, then solve for
the common multiplier.
\rev{This algebraic rearrangement of \eqref{eq:sharing-derivative} gives}
\begin{equation}
 \frac{\mathrm d\theta_j}{\mathrm d\alpha}
 =\frac{\displaystyle
 \frac{\mathrm dW_j}{\mathrm d\alpha}
 -\sum_{i\in{\cal C}}\frac{\mathrm dC_{ij}}{\mathrm d\alpha}
 -\theta_j\sum_{i\in{\cal F}}\frac{\mathrm dw_{ij}}{\mathrm d\alpha}}
 {\displaystyle\sum_{i\in{\cal F}}w_{ij}}.
 \label{eq:sharing-level-derivative}
\end{equation}
Subtract the capped grants and the response already supplied by changing
weights; normalize the remaining response by the free legs' total weight.
Evaluate the multiplier response, then each leg's response. If no leg is
uncapped, differentiate its caps directly; stable zero supply gives zero
rates. Differentiate \eqref{eq:lane-sharing} by the same balance with
equal lane weights. This supplies the current-rate channel of discharge.

Differentiate the service-amount sum \eqref{eq:service-amount} to turn
rate and duration changes into an actual transferred-mass response.
\rev{The next equation is its product-rule derivative:}
\begin{equation}
 \frac{\mathrm dy}{\mathrm d\alpha}
 =\sum_a\left(\delta_a\frac{\mathrm dr_a}{\mathrm d\alpha}
              +r_a\frac{\mathrm d\delta_a}{\mathrm d\alpha}\right).
 \label{eq:service-derivative}
\end{equation}
The first term changes rates over current durations; the second changes
time available at current rates.
\par\begin{samepage}
Equation \eqref{eq:head-time} selects
the duration branch. With the same controlling lane $\ell$ and stopping
condition,
\rev{differentiating that selected branch derives the three cases below:}
\begin{equation}
 \frac{\mathrm d\delta_a}{\mathrm d\alpha}
 =\begin{cases}
 \displaystyle\frac1{r_{\ell,a}}\frac{\mathrm dq_{\ell,a}}{\mathrm d\alpha}
 -\frac{q_{\ell,a}}{r_{\ell,a}^2}\frac{\mathrm dr_{\ell,a}}{\mathrm d\alpha},
 &\text{head clearance},\\[4pt]
 \displaystyle\frac1{r_{\ell,a}}\frac{\mathrm dS_{\ell,a}^{\rm rem}}{\mathrm d\alpha}
 -\frac{S_{\ell,a}^{\rm rem}}{r_{\ell,a}^2}\frac{\mathrm dr_{\ell,a}}{\mathrm d\alpha},
 &\text{sending allowance spent},\\[4pt]
 \displaystyle-\sum_{b<a}\frac{\mathrm d\delta_b}{\mathrm d\alpha},
 &\text{remaining timestep time}.
 \end{cases}
 \label{eq:duration-derivative}
\end{equation}
\par\end{samepage}
Head mass and remaining allowance here already include preceding
intervals' transfers. The first branch differentiates $r\delta=q$:
extra mass needs time, whereas faster service clears the head earlier.
It is the head-clearance path emphasized in the main report, not the
derivative of all minima at once. In the final branch,
$\sigma=\sum_{b<a}\delta_b$ and $\Delta t$ is fixed. Include any idle
remainder with zero rates. Tied stopping conditions retain their
directional/active-pattern treatment; the formula asserts no unique
ordinary derivative at a switch. These duration responses supply the
time channel in \eqref{eq:service-derivative}.

After service and admission, packet masses determine their cumulative
rank endpoints. Fixed bin boundaries give the following continuous path
through \rev{direct differentiation of the project representation \eqref{eq:grid}}:
\begin{equation}
 \begin{aligned}
 \frac{\mathrm dz^-_{\ell r}}{\mathrm d\alpha}
 &=\sum_{v<r}\frac{\mathrm dq_{\ell v}}{\mathrm d\alpha},\qquad
 \frac{\mathrm dz^+_{\ell r}}{\mathrm d\alpha}
 =\frac{\mathrm dz^-_{\ell r}}{\mathrm d\alpha}
                  +\frac{\mathrm dq_{\ell r}}{\mathrm d\alpha},\\
 \frac{\mathrm dM_{\ell cm}}{\mathrm d\alpha}
 &=\sum_{r:m_r=m}\frac{\mathrm d\omega_{\ell cr}}{\mathrm d\alpha},
 \qquad
 \frac{\mathrm dn_\ell}{\mathrm d\alpha}
 =\sum_{c,m}\frac{\mathrm dM_{\ell cm}}{\mathrm d\alpha}.
 \end{aligned}
 \label{eq:grid-derivative}
\end{equation}
Differentiate each overlap through its active min/max endpoints. The
cell boundary is fixed; the packet boundary moves. Fixed metadata are
not required to carry gradients. Exact-compatible merging preserves the
same overlap and timing map.

\par\begin{samepage}
For a positive-width readiness interval $[t_{\rm f},t_{\rm t}]$ and
fixed cutoff time $t$,
\rev{uniform readiness timing derives the interior fraction below; quotient and product rules then derive its sensitivities.}
The first line defines that fraction:
\begin{equation}
 \begin{aligned}
 f&=\frac{t-t_{\rm f}}{t_{\rm t}-t_{\rm f}},\\
 \frac{\mathrm df}{\mathrm d\alpha}
 &=\frac{\displaystyle
 -\frac{\mathrm dt_{\rm f}}{\mathrm d\alpha}
 -f\left(\frac{\mathrm dt_{\rm t}}{\mathrm d\alpha}
                 -\frac{\mathrm dt_{\rm f}}{\mathrm d\alpha}\right)}
 {t_{\rm t}-t_{\rm f}},\\
 \frac{\mathrm dq^{\rm ready}}{\mathrm d\alpha}
 &=f\frac{\mathrm dq}{\mathrm d\alpha}
                       +q\frac{\mathrm df}{\mathrm d\alpha}.
 \end{aligned}
 \label{eq:eligibility-derivative}
\end{equation}
\par\end{samepage}

Use whole/empty branches outside the interval. Point segments use their
boundary branch. For positive total entrance $b$, differentiate the
front-entry endpoint of \eqref{eq:entry-time}
\rev{by the quotient rule to obtain}

\begin{equation}
 \frac{\mathrm dt^{\rm in}_{\rm front}}{\mathrm d\alpha}
 =\Delta t\left(\frac1b\frac{\mathrm dz}{\mathrm d\alpha}
       -\frac z{b^2}\frac{\mathrm db}{\mathrm d\alpha}\right).
 \label{eq:entry-derivative}
\end{equation}

The tail replaces $z$ by $z+q^{\rm in}$. Zero entrance creates no segment.
The free-flow time is fixed. Delayed cumulative queries keep fixed
interpolation weights but differentiate stored boundary counts. On
positive eligible-mass and storage-limited receiving branches,
\rev{direct differentiation of the inherited \eqref{eq:sending} and \eqref{eq:receiving} gives:}
\begin{equation}
 \begin{aligned}
 \frac{\mathrm dA_{e,k}}{\mathrm d\alpha}
 &=\frac{\mathrm d{\cal H}^+_e(t_{k+1}-\tau^f_e)}{\mathrm d\alpha}
                     -\frac{\mathrm dN^-_{e,k}}{\mathrm d\alpha},\\
 \frac{\mathrm dR_{e,k}}{\mathrm d\alpha}
 &=\frac{\mathrm d{\cal H}^-_e(t_{k+1}-\tau^b_e)}{\mathrm d\alpha}
                     -\frac{\mathrm dN^+_{e,k}}{\mathrm d\alpha}.
 \end{aligned}
 \label{eq:ltm-derivative}
\end{equation}
Travel-limited sending follows $\mathrm dA/\mathrm d\alpha$; a fixed
capacity-limited sending or receiving branch has zero derivative.
Stable zero branches give zero availability response. These branch rules
are the derivatives of the existing LTM limits, not a new propagation law.
\par

\par\begin{samepage}
\section{Complete population time and continued state response}


Collect source masses, cumulative histories and continuous ledger
mass/timing fields in $X_k$. Per-lane grids are deterministic functions
of those data inside $F_k$. If $g_k$ denotes applied service allocation,
use $X_{k+1}=F_k(X_k,g_k)$ and the complete population account below.
\rev{The update map denotes the present model composition; the objective reuses the LTM project's source-plus-road time quadrature:}
\begin{equation}
 J_{\Delta t}=\Delta t\sum_{k=0}^{H-1}
 \left[\sum_o(Q_{o,k}+D_{o,k})+\sum_en_{e,k}\right].
 \label{eq:tst}
\end{equation}
$Q$ is existing source waiting, $D$ newly available release, $n$ all road
mass and $H$ the common complete-clearance horizon. The time charge has
vehicle-equivalent-second units. Source admission transfers ownership but
does not remove population from the account.
\par\end{samepage}

Let $c_k$ be the timestep charge in \eqref{eq:tst}. Differentiating
the displayed state update $X_{k+1}=F_k(X_k,g_k)$ and summing these charges
\rev{derives the following chain-rule identities from those state and objective definitions:}
\begin{equation}
 \begin{aligned}
 \frac{\mathrm dX_{k+1}}{\mathrm d\alpha}
 &=\frac{\partial F_k}{\partial X_k}
            \frac{\mathrm dX_k}{\mathrm d\alpha}
  +\frac{\partial F_k}{\partial g_k}
            \frac{\mathrm dg_k}{\mathrm d\alpha},\\
 \frac{\mathrm dJ}{\mathrm d\alpha}
 &=\sum_k\left(
 \frac{\partial c_k}{\partial X_k}\frac{\mathrm dX_k}{\mathrm d\alpha}
 +\frac{\partial c_k}{\partial g_k}\frac{\mathrm dg_k}{\mathrm d\alpha}
 \right).
 \end{aligned}
 \label{eq:state-sensitivity}
\end{equation}
During control $\mathrm dg_k/\mathrm d\alpha=d$; during fixed common
clearance it is zero, while the state response persists. Candidate-independent
empty loading/preparation gives initial response zero. At a fixed state
the present population charge has zero direct green partial. Stop-line
service responds directly; its later queue/history changes alter admission,
spillback, source waiting and eventual exit.

\par\begin{samepage}
For fixed releases, let $D_{\leq k}$ be cumulative available release and
$Z_k$ cumulative sink exits before update $k$. The conservation identity \eqref{eq:mass-identity}, after adding
update $k$'s new release, equates source-plus-road population in
\eqref{eq:tst} to $D_{\leq k}-Z_k$. With releases, timestep and common
horizon fixed,
\rev{the following objective expression and derivative are algebraic consequences of \eqref{eq:mass-identity} and \eqref{eq:tst}:}
\begin{equation}
 J_{\Delta t}=\Delta t\sum_{k=0}^{H-1}(D_{\leq k}-Z_k),\qquad
 \frac{\mathrm dJ_{\Delta t}}{\mathrm d\alpha}
 =-\Delta t\sum_{k=0}^{H-1}\frac{\mathrm dZ_k}{\mathrm d\alpha}.
 \label{eq:exit-objective}
\end{equation}
\par\end{samepage}

\par

Use a common integer horizon qualified to clear the compared action
neighborhood, followed by zero occupancy. Finite demand, acyclic routes,
positive continuation service and eventual sink availability support
qualification; no universal numerical horizon or clearance theorem is
asserted. A changing stop counter, residual cutoff or nonclearing case
cannot silently supply complete TST. The derivative domain includes
clearance and unchanged active signatures; switches need directional
branches, not a global smoothness claim.

\section{Matched comparisons and native qualification}

Models share population, geometry/permissions, action, objective,
preparation, continuation and a declared data protocol.
Original fine-grained series may pass through each model's own macro
converter/estimator. Declare retained variables, memory, update schedule
and calibration data. Native subsets are legitimate. Successful macro
inference counts against needing exact order; an exact ordered ledger
retained inside an estimator remains meso information.

Equal current totals do not establish equality of all lawful histories.
State-conditioned replay, observation-updated prediction and forecasting
from a common prefix are distinct tasks. State which observations are
fixed or regenerated for each derivative claim. A later observation
cannot count as a response predicted before it was observed.


The comparison models are aggregate link, structured movement-queue,
lane/storage aggregate and selective ordered-mass descriptions. They
explain the same physical ablations with different retained state.
Queue/storage explanations are expected to suffice when their mechanisms
are the only active effects; combined order/access cases test the extra
value of realized sequence.
\par

The movement-queue comparison uses
\href{https://arxiv.org/html/2310.12249v2}{Wei et al.'s link queue model},
with moving/queued state and native effective green. Its full causal
short-region/link/node law must be instantiated; the fixed-speed shortcut
is not assumed valid in every spillback regime.

The lane/storage comparison uses
\href{https://www.sciencedirect.com/science/article/pii/S0191261525000578}
{Gong et al.'s lane-exit allocation} with a sourced storage adaptation,
for which \href{https://attap.umd.edu/wp-content/uploads/2022/05/TRC522.pdf}
{Liu--Chang (2011)} is a candidate. Exit allocation is not finite-bay
entry, and native theorems do not automatically qualify that composition.
Exact implementations and adapters still require physical and task-specific
checks; the cited native model is not automatically a qualified baseline.

A V2 comparator failure must be understood before attributing a V3
advantage to retained order.

\section{Definition checks and external usefulness}

Local reference checks cover pool/feeder access, partial admission,
finite storage, current-head
sharing/reuse, route conservation and the source-counted objective.
V0 neutral cases need matched pure service classes and nonbinding or
equivalent receiving; one through/right green does not make differently
constrained exits identical. In that neutral domain the laws reduce to
the same sending/service/receiving minimum. Exact equivalence concerns
a matched aggregate backbone; a source-native propagation difference
must be qualified and cannot be credited to retained order.

Rotating the symmetric scenario and green groups must rotate results.
Equal lane totals do not justify one larger FIFO lane without
corresponding order, readiness, access, service and receiving.
Hand/reference checks establish definition correctness, not realism.

For the richer-reference comparison, use one fixed executable mapping. With fixed cycle
$C^{\rm cyc}$, lost intervals totaling $L<C^{\rm cyc}$, order and offset,
\rev{the conversion below is derived from the task's fixed usable-time budget and green-fraction definition:}
\begin{equation}
 G_p=(C^{\rm cyc}-L)g_p,\qquad
 s^{\rm eff}_{\ell m}=s^{\rm phys}_{\ell m}
                         (C^{\rm cyc}-L)/C^{\rm cyc}.
 \label{eq:reference-action}
\end{equation}
Declare program admissibility and the common switch to fixed clearance;
do not silently insert minimum green to repair an API.
This matches nominal average opportunity, not microscopic startup/access
or every phase effect. Official
\href{https://sumo.dlr.de/docs/Simulation/Traffic_Lights.html}
{SUMO signal definitions} support the interface, not empirical equivalence.
\rev{Those definitions describe how to encode the program; they are not the source of the budget identities above.}

Use a finite master route population and common desired releases.
Integer reference vehicles are exogenous population, not rounded internal
mass. Fix native lane behavior independently of candidate directions;
do not replay the ordered model's private assignments into only one model.
Use matched predeclared reference seeds and held-out family IDs,
keeping order pairs together.

Reference time is desired-release-to-arrival time:
\texttt{duration} plus \texttt{departDelay}, from official
\href{https://sumo.dlr.de/docs/Simulation/Output/TripInfo.html}{TripInfo}.
Waiting is already inside duration. Retain uninserted/unfinished members
in master-population accounting; clear all demand or report noncompletion,
never a completed-trip-only objective.

For each feasible direction, compare predicted sign/ranking and scale
with actual richer-evaluator complete-time change under the same action
mapping. Disclose observation conditioning. Only after forward/reference
and derivative checks pass should matched optimization and held-out
usefulness be tested.

\section{Scope and evidence limits}

The model specifies eligible packet inspection, movement-specific waiting
priority, actual-admission balancing, physical lane access, elapsed
service, interpolated histories and complete population-time sensitivity.
Its parallel approach lanes retain distinct occupancy, storage and order.
The bay feeder can obstruct one part of the approach while another lane
remains usable, without giving resident traffic an unmodeled bypass.

Feasibility and unique-selection arguments concern the mathematical
operators. The analytic examples illustrate local consequences.
Implementation correctness, numerical derivative agreement, qualified
baselines and external control usefulness require separate evidence.

\end{document}
